\chapter{Letters and Proposals}\label{app:g}

\section{Assistance Letter Request for the PPA}
\label{appendix-g-assistance-letter}

\noindent January 22, 2026

\noindent\textbf{ENGR. KRIS LLOYD T. VILLAMONTE}\\
\emph{Director}, Physical Plant Administration (PPA)\\
Ateneo de Naga University\\
Naga City 4400 Philippines

\noindent Dear Engr. Villamonte,

Greetings of Peace!

We are writing to you upon the specific recommendation of our thesis panel to
request the technical assistance of the Physical Plant Administration (PPA) for
our thesis study, \emph{TinyML Smart Plug for Complementary Protection Against
Progressive Fire-Inducing Electrical Faults in Residential Circuits}.

Our research aims to develop a device that detects electrical faults that
commonly cause fire, specifically series arc faults, socket heating, and
overload that standard breakers often miss. To acquire data for machine
learning model training and to validate the performance of the resulting
device, we must simulate hazardous conditions, including generating arc faults
and inducing overload and heat in the electrical system.

Given the risks involved, we need the support of personnel who are skilled in
electrical safety to execute, control, and supervise the simulation of
electrical faults.

Attached is our Proposed Safety Guidelines and Testing Protocols, which
details the specific dates, times, and locations for the testing and our
general plan in conducting the testing.

We would be deeply grateful for your support in safely finishing our thesis
study.

\noindent Respectfully yours,

\vspace{4\baselineskip}
\noindent\textbf{FRANCIS M. PE\~{N}A}\\
\emph{4 - BS ECE Student}

\vspace{2\baselineskip}
\noindent\textbf{NOTED BY:}

\vspace{3\baselineskip}
\noindent
\begin{minipage}[t]{0.47\textwidth}
\centering
\textbf{DR. REFERENDO D. SORIANO}\\
\emph{Thesis Adviser}\\
Date Signed: \rule{1.45in}{0.4pt}
\end{minipage}\hfill
\begin{minipage}[t]{0.47\textwidth}
\centering
\textbf{DR. CELESTE M. OJEDA}\\
\emph{Chairperson}, ECE/CpE Department\\
Date Signed: \rule{1.45in}{0.4pt}
\end{minipage}

\section{Safety Guideline Proposal to the PPA}
\label{appendix-g-safety-guideline-proposal}

\begin{center}
\textbf{PROPOSED SAFETY GUIDELINES AND TESTING PROTOCOLS}\\
For the Conduct of Electrical Fault Simulation during Data Collection and Testing
\end{center}

\subsection*{I. Title and Proponents}

The research project, titled \emph{TinyML Smart Plug for Complementary
Protection Against Progressive Fire-Inducing Electrical Faults in Residential
Circuits}, is presented by \textbf{Alex A. Alferez, Jett M. Lumabi, and
Francis M. Pe\~{n}a}, fourth-year BS Electronics Engineering students from the
Department of Electronics and Computer Engineering at Ateneo de Naga
University. Conducted under the advisement of Dr. Referendo D. Soriano, the
study is scheduled for its primary development and validation phase from
\textbf{January 2026 to May 2026}.

\subsection*{II. Overview}

The study aims to develop an intelligent, ESP32-based smart plug capable of
detecting fire-inducing electrical faults that standard circuit breakers often
miss. These faults are series arc faults, silent sustained overloads, and
contact heating.

Since these faults are inherently dangerous, this proposal seeks the assistance
of trained PPA personnel who are skilled in safely handling electrical circuits
to ensure the safe and controlled execution of laboratory simulations required
for data collection and system validation.

\subsection*{III. Rationale}

Validating a solution for these electrical faults requires generating them in a
controlled environment. However, creating actual arcs, inducing overloads, and
artificially degrading electrical sockets involve significant risks including
fire, shock, and equipment damage. To mitigate potential risks and protect
students, laboratory equipment, and facilities, all high-voltage procedures and
fault simulations will be conducted under the direct supervision and execution
of authorized PPA personnel.

\subsection*{IV. Objectives}

The objective of this proposal is to safely collect actual electrical waveform
data of hazardous faults to train and test the TinyML fault-detection system,
specifically:
\begin{enumerate}
\item To simulate sustained overloads from 10A to 15A to test overload
behavior, current-sensor accuracy, and protective response;
\item To induce contact heating above 70\degreeC on test terminals or a
degraded socket to validate thermal sensing and estimation;
\item To simulate series arc faults using a controlled arc generator setup
under PPA supervision; and
\item To establish a safe protocol where PPA personnel manage live wires and
fault simulations while students handle data logging.
\end{enumerate}

\subsection*{V. Scope and Limitations}

\begin{enumerate}
\item The request covers the period from January 2026 to May 2026.
\item PPA assistance is requested specifically for live fault generation,
emergency switching, and setup inspection.
\item Testing is limited to single-phase 230V AC in a controlled laboratory
environment.
\item PPA staff are not expected to debug ESP32 code or electronics; their role
is to manage the AC mains supply and fault-generation rigs for safety.
\end{enumerate}

\subsection*{VI. Methodology and Safety Protocols}

The following protocols will be observed throughout the testing period:
\begin{enumerate}
\item No high-voltage testing or fault simulation will commence without an
authorized PPA staff member present.
\item The testing setup will use an accessible circuit breaker, and the D315
workshop breaker must remain accessible to laboratory custodians and PPA
personnel.
\item All personnel inside the testing perimeter must wear rubber-soled safety
shoes, insulated gloves, and safety glasses.
\item Students will handle data logging and software monitoring. PPA staff will
handle live connections and fault triggers.
\end{enumerate}

\begin{table}[H]
  \centering
  \caption{Hazard identification and mitigation strategy}
  \label{tab:app-g-hazard-mitigation}
  \begin{tabular}{p{0.18\textwidth}p{0.30\textwidth}p{0.42\textwidth}}
    \toprule
    \textbf{Fault type} & \textbf{Identified hazard} & \textbf{Mitigation strategy} \\
    \midrule
    Overload & Wire overheating and potential fire & Use a controlled load bank or source, maintain real-time thermal monitoring, and keep breaker access available. \\
    Arc fault & Plasma burns, UV radiation, and fire ignition & Require safety goggles and insulated gloves, keep a fire extinguisher on standby, and use an arc containment cover or box. \\
    Contact heating & Burns and fumes from heated or melting materials & Provide ventilation, use non-contact temperature checks, use face masks as needed, and follow hot-work precautions. \\
    \bottomrule
  \end{tabular}
\end{table}

\subsection*{VII. Proposed Schedule}

The following timeline details the activities requiring PPA presence. Dates are
subject to personnel availability, institutional announcements, weather
interruptions, and holidays.
\begin{table}[H]
  \centering
  \caption{Proposed schedule of activities requiring PPA assistance}
  \label{tab:app-g-proposed-schedule}
  \scriptsize
  \setlength{\tabcolsep}{3pt}
  \begin{tabular}{p{0.20\textwidth}p{0.15\textwidth}p{0.15\textwidth}p{0.34\textwidth}}
    \toprule
    \textbf{Activity} & \textbf{Dates} & \textbf{Location} & \textbf{PPA involvement} \\
    \midrule
    Checking the testing setup & February 26--27, 2026 & D315 Workshop & Confirm the safety of the researchers' testing setup. \\
    Data collection for overload & March 2--6, 2026 & D315 Workshop & Assist in simulating overload using appliances or a high-power load while breaker access is maintained. \\
    Data collection for contact heating & March 9--13, 2026 & D315 Workshop & Assist in simulating contact heating through a degraded socket or controlled heating of contacts. \\
    Data collection for arc fault & March 16--April 1, 2026 & D315 Workshop & Assist in simulating a series arc fault by operating the covered arc-fault generator. \\
    Model and device testing & April 6--20, 2026 & D315 Workshop and Electrical Warehouse & Assist during induced-fault tests used to evaluate the device algorithm. \\
    \bottomrule
  \end{tabular}
\end{table}
\begin{table}[H]
  \centering
  \caption{Proposed daily testing hours}
  \label{tab:app-g-testing-hours}
  \begin{tabular}{ccccc}
    \toprule
    \textbf{Monday} & \textbf{Tuesday} & \textbf{Wednesday} & \textbf{Thursday} & \textbf{Friday} \\
    \midrule
    9AM--3PM & 9AM--5PM & 9AM--3PM & 9AM--5PM & 9AM--3PM \\
    \bottomrule
  \end{tabular}
\end{table}

\vspace{1\baselineskip}
\noindent\textbf{NOTED BY:}

\vspace{2\baselineskip}
\noindent
\begin{minipage}[t]{0.47\textwidth}
\centering
\textbf{DR. REFERENDO D. SORIANO}\\
\emph{Thesis Adviser}\\
Date Signed: \rule{1.45in}{0.4pt}
\end{minipage}\hfill
\begin{minipage}[t]{0.47\textwidth}
\centering
\textbf{DR. CELESTE M. OJEDA}\\
\emph{Chairperson}, ECE/CpE Department\\
Date Signed: \rule{1.45in}{0.4pt}
\end{minipage}

\vspace{1\baselineskip}
\noindent\textbf{APPROVED BY:}

\vspace{2\baselineskip}
\noindent
\begin{minipage}[t]{0.58\textwidth}
\centering
\textbf{ENGR. KRIS LLOYD T. VILLAMONTE}\\
\emph{Director}, Physical Plant Administration (PPA)\\
Date Signed: \rule{1.45in}{0.4pt}
\end{minipage}
